% Cover letter and point-by-point reviewer responses.
% Reviewer text is verbatim from paper_review.txt; \reply{...} blocks
% are scaffolded as TODO placeholders to be filled in during the revision.

\begin{minipage}[b]{2.5in}
  Resubmission Cover Letter \\
  {\it Bioinformatics}
\end{minipage}
\hfill
\begin{minipage}[b]{5.5in}
  Andrew Kern (on behalf of all authors) \\
  \today
\end{minipage}

\vskip 2em

\noindent
{\bf To the Editor(s) -- }

\vskip 1em

We are writing to submit a revised version of our manuscript titled
``MKado: a toolkit for McDonald-Kreitman tests of natural selection''.
We thank the editor and both reviewers for their thoughtful and
constructive comments, which have substantially strengthened the
manuscript and the underlying software package. The most substantive
changes in this revision are: (i) a new aggregated-genes asymptotic-MK
runtime benchmark addressing Reviewer~1, point~5; (ii) implementation of
the Uricchio et al.~(2019) ``above-$x$'' cumulative SFS for the
asymptotic test, addressing Reviewer~1, point~4; (iii) a polyDFE-style
case-resampling bootstrap CI for asymptotic alpha, addressing the second
half of Reviewer~1, point~5; and (iv) prominent VCF-input and
allele-polarization documentation in both manuscript and docs,
addressing Reviewer~1, points~1--2 and Reviewer~2.

In addition to a Response to Reviewers (in which page/line numbers refer
to the revised manuscript file), we provide a marked-up PDF showing the
differences relative to the original submission.

\vskip 2em

\noindent\hspace{4em}
\begin{minipage}{4.5in}
\noindent
{\bf Sincerely,}

\vskip 2em

{\bf
  Andrew Kern (on behalf of all authors)
}\\
\end{minipage}

\vskip 4em

\pagebreak

%%%%%%%%%%%%%%%%%%%%%%%%%%%%%%%%%%%%%%%%%%%%%%%%%%%%%%%%%%%%%%%%%%%%%%%%%
\reviewersection{1}

\begin{itquote}
McDonald and Kreitman (MK) test is one of the most widely used methods
to test for natural selection. Since its original publication, several
studies have proven its potential for detecting positive selection.
Nowadays, we have several extensions that mitigate the main issue of
the MK test: the presence of slightly deleterious mutations. Here, the
authors present a fast, user-friendly, and well-documented software to
automatize MK tests. The software not only includes multiple MK tests
extensions per gene or aggregated datasets, but also data processing,
multiple testing correction, and visualization features. However, I
have a few minor concerns that I hope the authors can address. Most of
them should be straightforward to test or implement.
\end{itquote}

\reply{We thank the reviewer for the positive assessment and for the
detailed, constructive comments. Each point is addressed in turn below;
the substantive items (especially points~4 and~5) drove the
two largest additions in this revision (cumulative-SFS option and a
new aggregated-genes runtime benchmark with bootstrap CIs).}

%---------- R1.1 ----------
\begin{point}{} % 1.1 -- VCF support visible in docs but not in manuscript
As the authors claim, most of the current software is fragmented. Some
tools can deal with VCF or multi-FASTA data (\texttt{degennotate},
\texttt{fastDFE}, \texttt{iMKT}; \url{https://doi.org/10.1093/molbev/msad270},
\url{https://doi.org/10.1093/molbev/msae070},
\url{https://doi.org/10.1093/nar/gkz372}), some have already implemented
several of the MK tests presented here (\texttt{MKtest.jl};
\url{https://doi.org/10.1093/g3journal/jkae031}), others can polarise
alleles (\texttt{fastDFE}; \url{https://doi.org/10.1093/molbev/msae070}).
\mkado unifies it on a user-friendly Python package. The documentation
shows that \mkado can deal with VCF files similarly to \texttt{fastDFE},
but I cannot see this mentioned in the paper. The authors should
include it in the main text, since handling both VCF and multi-FASTA
files is an important feature.
\end{point}

\reply{TODO: describe added VCF paragraph in Features section, new row
in the comparison table, and citations to the named alternatives.}

%---------- R1.2 ----------
\begin{point}{} % 1.2 -- automatic allele polarization vs polarized MK
In addition to describing the Polarised MK Test, the authors should
state in the main text that \mkado uses outgroup sequences to polarize
alleles. As far as I understand, this is done automatically and is
well described in the software documentation. Nonetheless, I would
recommend that the authors clarify the automatic allele polarization
required to run tests such as the asymptotic or imputed MK tests
(where polarization is mandatory or recommended), as distinct from the
polarised MK test, which specifically tests for selection along a
lineage. As I said, this is already well explained in the
documentation, so it should be straightforward to include in the paper
to avoid any misunderstanding.
\end{point}

\reply{TODO: describe new ``Allele polarization'' subsection in Methods.}

%---------- R1.3 ----------
\begin{point}{} % 1.3 -- TG alpha + imputed / FWW combination
When using the Tarone-Greenland Alpha, is the standard MK test
applied? Would it not be beneficial to correct for slightly deleterious
mutations by simply using the imputed or Fay, Wyckoff and Wu frequency
threshold correction?
\end{point}

\reply{TODO: explain that \texttt{--alpha-tg --min-freq X} already
applies the FWW frequency-threshold correction within
$\alpha_\text{TG}$; decline the imputed $\times$ TG combination as an
unvalidated novel estimator.}

%---------- R1.4 ----------
\begin{point}{} % 1.4 -- Uricchio et al. 2019 above-x SFS option
The authors should include an option to modify the SFS as described in
Uricchio et al.\ (2019) (see supplementary material,
\url{https://doi.org/10.1038/s41559-019-0890-6}):

\begin{itquote}
We note that the original asymptotic-MK approach takes $P_N(x)$ and
$P_S(x)$ as the number of polymorphic sites at frequency $x$ rather
than above $x$, but this approach scales poorly as sample size
increases since most common allele frequencies $x$ have very few
polymorphic sites in large samples. We therefore define $P_N(x)$ and
$P_S(x)$ as stated above since these quantities trivially have the
same asymptote but are less affected by changing sample size.
\end{itquote}

The asymptotic estimation would greatly benefit from this modification,
producing more robust results. The authors can follow
\url{https://github.com/jmurga/MKtest.jl/blob/main/src/polymorphism.jl}
to implement it.
\end{point}

\reply{TODO: describe new \texttt{--sfs-mode}\,\{\texttt{at},\texttt{above}\}
flag, default \texttt{at} (Messer \& Petrov 2013) for backward
compatibility, with the cumulative inclusive-tail variant available
under \texttt{above}; cite Uricchio 2019.}

%---------- R1.5 ----------
\begin{point}{} % 1.5 -- runtime benchmarks: composition + aggregated
Do the runtime plots show only the batch processing of the 12,437
human genes, or also the execution of the imputed or standard MK test?
If the tests are not included on runtime, I would recommend that the
authors incorporate both processing and MK test estimation, since this
should not substantially affect the benchmark. Moreover, I'm assuming
the author are not providing runtime benchmark for asymptotic
estimations. Because asymptotic estimation would probably fail at the
gene level or provide unreliable results (as the authors properly
note), I will recommend two approaches:

\begin{itemize}
\item[(i)] Following Murga-Moreno et al.\ (2022)
(\url{https://doi.org/10.1093/g3journal/jkac206}), the authors can
aggregate at least 1{,}000 random genes to ensure reliable estimation
on human data and run the benchmark over a reasonable number of
replicates.
\item[(ii)] They can add the polyDFE bootstrapping strategy to resample
the SFS over the 12{,}437 aggregated genes. This is a common strategy
for obtaining reliable alpha confidence interval (CI) estimates (e.g.,
Castellano et al.\ 2019;
\url{https://doi.org/10.1534/genetics.119.302494}). The authors
currently estimate CI via Monte Carlo, but adding a bootstrap
resampling strategy would benefit both runtime benchmarks and the
reliability of alpha estimates for any MK test provided. See the
\texttt{bootstrapData} function at
\url{https://github.com/paula-tataru/polyDFE/blob/master/postprocessing.R}
or
\url{https://fastdfe.readthedocs.io/en/latest/_modules/fastdfe/spectrum.html#Spectrum.resample}.
\end{itemize}

This way, non-expert users would not face runtimes longer than those
shown in Figure S1. Since the software scaling is nearly perfect,
testing approximately 1{,}000 aggregated datasets should be
informative enough.
\end{point}

\reply{TODO: describe (a) clarified composition of existing benchmarks
(parsing $+$ site counting $+$ test $+$ output), (b) new
aggregated-asymptotic benchmark and supplementary figure built on
the 1{,}000-gene-per-replicate regime
(\texttt{examples/benchmarks/aggregated\_asymptotic\_runtime.py}), and
(c) new \texttt{--ci-method bootstrap} option for asymptotic alpha,
implementing case-resampling refit per replicate.}

%---------- R1.6 ----------
\begin{point}{} % 1.6 -- Ls, Ln; omega_a, omega_na decomposition
Since the authors already manually estimate the shortest mutational
paths and weight synonymous and nonsynonymous sites following Nei and
Gojobori (1986), can they include the total number of synonymous and
nonsynonymous sites ($L_s$, $L_n$)? It should then be straightforward
to also implement $\omega_a$ and $\omega_{na}$, as decomposed in
Coronado-Zamora et al.\ (2019)
(\url{https://academic.oup.com/gbe/article/11/5/1463/5480667}):
\[
\omega_a = \alpha \cdot \omega, \qquad
\omega_{na} = (1 - \alpha) \cdot \omega.
\]
\end{point}

\reply{TODO: describe new \texttt{omega} command computing $L_n$, $L_s$,
$\omega = (D_n / L_n) / (D_s / L_s)$ and the
$\omega_a$ / $\omega_{na}$ decomposition with confidence intervals.}

%---------- R1.7 ----------
\begin{point}{} % 1.7 -- OrthoMaM polymorphism projection writeup
The authors followed an intelligent strategy using the OrthoMaM
database to obtain human $D_n$, $D_s$, $P_n$, and $P_s$ counts. I
think projecting the polymorphic data onto the ortholog alignments
could provide an useful feature when working with non-model species
available within OrthoMaM database. I understand that including such
a feature will be challenging, but a more detailed explanation of the
polymorphism projection procedure would help users working with
non-model organisms.
\end{point}

\reply{TODO: point at new ``Working with ortholog alignments'' section
in the documentation tutorial and a methods paragraph that summarizes
the OrthoMaM coordinate-mapping pipeline (alignment source, VCF
liftover, indel handling, multi-transcript reconciliation). Note that
the standalone \texttt{mkado project} subcommand is deferred to a
future release per the reviewer's allowance.}

%%%%%%%%%%%%%%%%%%%%%%%%%%%%%%%%%%%%%%%%%%%%%%%%%%%%%%%%%%%%%%%%%%%%%%%%%
\reviewersection{2}

\begin{itquote}
This is a concise manuscript reporting a standalone package
specifically designed for computing several flavours of the classical
McDonald-Kreitman test. The length of the manuscript, level of details
and information provided are appropriate. I agree with the authors
that such as tool can be a benefit for the community by filling a gap
in the set of available programs. I would say that \mkado fulfils its
tasks successfully. I had the chance to test the package and I found
it satisfying so my review will be positive. Nevertheless I will seize
this opportunity to provide feedback after looking at the code and
testing the tool.
\end{itquote}

\reply{We thank the reviewer for the positive overall assessment and
for testing the package end-to-end. The detailed, code-level feedback
was particularly valuable; each point is addressed below.}

%---------- R2.1 ----------
\begin{point}{} % 2.1 -- formal definition of Pn and Ps
[The text has a good structure.] However, I would recommend that both
the article and the documentation describe formally the basic
statistics used. There are detailed formulas for variants of the MK
tests, $D_n$ and $D_s$ are clearly defined, but $P_n$ and $P_s$ are
not, actually (or not completely).
\end{point}

\reply{TODO: describe new explicit definitions of $P_n$ and $P_s$ in
the manuscript Methods (with handling of singletons, the
\texttt{--pool-polymorphisms} option, and the derived-allele-frequency
threshold) and the matching tightening of dataclass docstrings in the
package.}

%---------- R2.2 ----------
\begin{point}{} % 2.2 -- VCF support not mentioned in manuscript
I am surprised that the possibility of importing VCF data is not
mentioned in the manuscript. I suppose that this was a recent addition
to the package. This would be an attractive addition to the article.
\end{point}

\reply{TODO: cross-reference response to Reviewer~1, point~1 (added
VCF paragraph and comparison-table row).}

%---------- R2.3 ----------
\begin{point}{} % 2.3 -- Figure 1 caption: two-panel
Figure 1 should be cited slightly differently as it shows both the
volcano plot and the asymptotic MK test display.
\end{point}

\reply{TODO: describe the caption rewrite (left/right panel
description) and the in-text references that no longer mistakenly
refer to a single panel.}

%---------- R2.4 ----------
\begin{point}{} % 2.4 -- __init__.py outdated (version, package name)
The main \texttt{\_\_init\_\_.py} file is outdated as of this writing.
First, the hardcoded version number is still 0.3.0 while the package
has been ported to 0.4.0. Second, the package name in the file
docstring is still ``mikado'', which is the original name of the
project.
\end{point}

\reply{TODO: confirm the version-string and docstring fix in
\texttt{src/mkado/\_\_init\_\_.py} (single-source-of-truth via
\texttt{importlib.metadata.version}).}

%---------- R2.5 ----------
\begin{point}{} % 2.5 -- shell completion options unclear; slow
I noticed that the umbrella command has options to manage shell
completion. They were not clear to me at first sight and it is rather
unexpected to find them as the option of the umbrella command, which
is not expected to do anything. Maybe a ``setup'' command would
clarify it but this is not essential. In any case I recommend adding a
sentence to the help page to explain that these options allow to
configure automatic shell completion. On my system, shell completion
was painfully slow so the feature was doing more harm than good.
\end{point}

\reply{TODO: describe the help-text addition for
\texttt{--install-completion} / \texttt{--show-completion} on the
umbrella command, plus a new ``Shell completion'' subsection in the
docs noting that completion may be slow on some shells and can simply
be left uninstalled.}

%---------- R2.6 ----------
\begin{point}{} % 2.6 -- subcommand-without-args behaviour
When called without argument, the different commands end up with an
error, which is somewhat inconsistent with the behaviour of the
umbrella command (which displays the help page). But it is a perfectly
acceptable behaviour and up to the authors. Commands have logical and
relevant options.
\end{point}

\reply{TODO: short response noting the reviewer's allowance to keep
the existing behaviour, or describe a small CLI tweak that prints help
when a required positional is missing.}

%---------- R2.7 ----------
\begin{point}{} % 2.7 -- output filename specification
The commands write a progress bar and some feedback to the standard
error stream, obviously to enable standard output redirection to
capture the output. Unless I missed it, there is no way to specify
[an] output file name. The choice that has been made is not perfectly
complying with conventions regarding the use of the standard error
stream but it is not unusual either.
\end{point}

\reply{TODO: describe the new \texttt{--output} / \texttt{-o} flag (or
explicitly note that stdout redirection plus the existing
\texttt{--format} flag already covers this case, if we decide not to
add a new flag).}

%---------- R2.8 ----------
\begin{point}{} % 2.8 -- ruff out of user manual; "third outgroup"
The package is adequately documented by a fairly complete README
available on the git and PyPI repositories, and, more in more details
by a readthedocs page. [A] small remark would be that I don't think
it is relevant to explain how to use ruff in the manual, as this is
not needed to use \mkado. By contrast, the shell completion options
could be mentioned. There is somewhere the term ``third outgroup''
which I think is incorrect. Apart for this, the documentation is
clear and complete, with very helpful commands.
\end{point}

\reply{TODO: confirm ruff usage moved out of the user-facing
installation page (now in CONTRIBUTING.md), shell completion documented
(see R2.5), and ``third outgroup'' corrected to ``second outgroup'' in
\texttt{docs/index.rst}.}

%---------- R2.9 ----------
\begin{point}{} % 2.9 -- aggregate option meaning + range/regex matchers
Only the meaning of the ``aggregate'' option was not clear to me. As
samples are often grouped by population, it could be a possibility to
allow specifying the ingroup and outgroup by range indexes, in top of
the two possibilities already provided. This might be a feature
request. Another suggestion is to support regular expressions to
identify ingroup and outgroup samples in addition to an exact
substring match.
\end{point}

\reply{TODO: (a) describe the new tutorial subsection that explicitly
contrasts \texttt{--aggregate} and \texttt{--per-gene}; (b) describe
the new index-range matchers; (c) decline regex matchers with the
substring-vs-regex trade-off rationale.}

%---------- R2.10 ----------
\begin{point}{} % 2.10 -- Figure S1 linear scaling claim
For large datasets it is still needed to use parallelization. However,
Python isn't (currently) really good at parallelization. My feeling is
that processing one alignment doesn't involve a lot of computation
compared with the cost of spawning a worker (which entails creating a
whole Python process) so the number of workers should be orders of
magnitude below the number of genes (to avoid a waste of resources for
little benefit). Figure S1 actually shows that the relationship is all
but linear so I am not sure to agree with the sentence claiming that
speedup is linear until 32 workers.
\end{point}

\reply{TODO: describe the rewritten parallelism paragraph (now
quantitative: ``approximately $N\times$ speedup at 16 workers, with
diminishing returns past 32 due to per-process spawn overhead'') and
the corresponding rewrite of the Figure S1 caption.}
